% --- ОБЩИЕ НАСТРОЙКИ ТЕКСТА ---
\onehalfspacing
\setlength{\parindent}{1.25cm} 

\sloppy 
\hyphenpenalty=1000 
\tolerance=2000

% --- ФИКС РАЗМЕРОВ ЗАГОЛОВКОВ (14pt везде) ---
\setkomafont{disposition}{\rmfamily\bfseries\normalsize}
\setkomafont{chapter}{\rmfamily\bfseries\normalsize}
\setkomafont{section}{\rmfamily\bfseries\normalsize}
\setkomafont{subsection}{\rmfamily\bfseries\normalsize}

% Настройка оглавления (текст обычный, заголовок жирный)
\setkomafont{chapterentry}{\normalfont\normalsize}
\setkomafont{chapterentrypagenumber}{\normalfont\normalsize}
\KOMAoptions{toc=chapterentrydotfill}

% --- НАСТРОЙКА ОГЛАВЛЕНИЯ (СОДЕРЖАНИЕ) ---
\renewcaptionname{russian}{\contentsname}{\bfseries СОДЕРЖАНИЕ}
\BeforeTOCHead[toc]{\renewcommand*{\raggedsection}{\centering}}

% --- МАГИЯ ОТСТУПОВ ЗАГОЛОВКОВ (1.25 см) ---
% Эта команда перехватывает отрисовку заголовка и вставляет отступ
\makeatletter
\renewcommand{\sectionlinesformat}[4]{%
  \ifstr{#1}{chapter}{%
    % Для глав (если они не центрированы)
    \hspace{1.25cm}#3#4%
  }{%
    % Для разделов и подразделов
    \hspace{1.25cm}#3#4%
  }%
}
\makeatother

% Убираем точку после номера (1.1 вместо 1.1.)
\KOMAoptions{numbers=noendperiod}

% Базовые интервалы для заголовков
\RedeclareSectionCommand[beforeskip=12pt, afterskip=6pt]{chapter}
\RedeclareSectionCommand[beforeskip=12pt, afterskip=6pt]{section}
\RedeclareSectionCommand[beforeskip=12pt, afterskip=6pt]{subsection}

% --- СПИСКИ (enumitem) ---
\setlist{
	leftmargin=0pt,
	nosep,
	wide=\parindent,
	labelsep=0.5em,
	listparindent=\parindent
}
\setlist[itemize,1]{label=--} 
\setlist[enumerate,1]{label=\arabic*.}

% --- ПОДПИСИ И ТАБЛИЦЫ ---
\captionsetup{justification=centering, singlelinecheck=false}
\DeclareCaptionLabelFormat{simple}{#1~#2}
\captionsetup{labelformat=simple}
\DeclareCaptionLabelSeparator{emdash}{ --- }
\captionsetup[table]{
	labelsep=emdash,
	justificat